\chapter{Design}
\label{chap:design}

This chapter sets out the experimental design: the architecture under test, the
partitioning schemes, the evaluation protocol, and the reasoning behind each
choice.

\section{Architecture under test}
\label{sec:architecture}

PerMFL holds three models, each a complete copy of the network rather than a
subset of its layers. A device model $\theta_{i,j}$ belongs to device $j$ in
team $i$, a team model $w_i$ to each of the $M$ teams, and a global model $x$
to the server. Figure~\ref{fig:tiers} shows the arrangement and the direction
of each coupling.

The training loop nests three counters. Each global round $t$ runs $K$ team
rounds, and each team round runs $L$ local gradient steps on every
participating device. At the start of a global round every team model is
overwritten with the current global model. At the start of every team round
every device model is overwritten with its team's model. A device therefore
travels at most $L$ steps from its team model before being reset, a property
that Section~\ref{sec:depth} shows to be the dominant influence on personalised
performance.

Communication cost follows from the loop structure. Each team round broadcasts
the team model to its devices and collects their updates, so device-to-server
transfers per global round are $2KN$ for $N$ devices in total. The number of
teams cancels: five teams of four and one team of twenty cost the same.

\section{Partitioning schemes}
\label{sec:partitioning}

How clients are constructed determines what any clustering method can recover
and how much a group-level model can contribute. Four schemes are implemented.

\subsection{Attack domain}

Each attack family is assigned to a group of clients, following CFMD-i's first
scenario. Benign traffic is distributed uniformly across every client, since a
detector with no notion of normal traffic has no decision boundary to learn.
Clients are dealt to domains in rotation and then shuffled, so that client index
does not encode domain membership. Without the shuffle, a clustering that
recovered index order would be indistinguishable from one that recovered
structure.

This scheme is the only one that supplies ground truth. Every client carries a
known domain label, so a derived partition can be scored against it with the
adjusted Rand index rather than argued about through downstream accuracy.

\subsection{Label skew}

Each client receives a fixed number $k$ of attack classes, dealt in rotation so
that neighbouring clients overlap and every class has a similar number of
holders. Sweeping $k$ produces a heterogeneity ladder, following the approach in
Stones From Other Hills, which reports a federated baseline, a local baseline
and its own method at each level.

\subsection{Dirichlet}

The standard approach in the clustered federated learning literature. Class
proportions per client are drawn from a Dirichlet distribution with
concentration $\alpha$, lower values producing more skew. It is included for
comparability, with the caveat noted in Section~\ref{sec:gap}: it produces a
continuum rather than discrete groups, so it cannot support a ground-truth
clustering evaluation.

\subsection{Host based}

One client per destination host, which is what a deployed system looks like:
one agent per machine. Most victim hosts in these captures see a single attack
type, so hosts are selected by the entropy of their label distribution, an
approach taken from DP-FL. Without that filter every client would hold one
class and the partition would be degenerate.

\section{Controlling and measuring heterogeneity}
\label{sec:hetero-control}

Two independent controls are provided.

The first is the partition itself, through $k$ or $\alpha$. The second is the
proportion of benign traffic retained per client. At eighty-two per cent benign,
every client's label distribution is dominated by the same class, which
compresses any divergence measure regardless of how the attacks are split.
Thinning benign traffic on the training split makes the underlying structure
visible without altering the partition. It is applied to training data only:
the evaluation set must reflect operational traffic, and thinning it would
inflate macro F1 by shrinking the majority class in the denominator.

Heterogeneity is reported as the mean pairwise Jensen-Shannon divergence
between client label distributions, which is zero when clients are identical
and one when they are disjoint, alongside the effective number of classes per
client.

\section{Evaluation protocol}
\label{sec:eval-protocol}

\subsection{Metrics}

Accuracy is reported but is not the headline. On a corpus that is eighty-two
per cent benign, a classifier that never raises an alert scores 0.8171. The
reason is structural: accuracy uses the true negative cell, which dominates
under imbalance. For a single attack class in the CICIDS2017 evaluation set, a
detector that never fires records per-class accuracy between 0.955 and 0.992
while recording zero recall.

Macro-averaged F1 is the headline metric. It weights every class equally and,
uniquely among the composite measures considered, does not use the true
negative cell at all. Recall is reported as the detection rate an operator
tunes toward, and false positive rate as the alert volume they tune against.

Every result is reported against two anchors: the floor a majority-class
classifier achieves, and the ceiling the same model and optimiser reach when
trained centrally on the pooled data.

\subsection{Splits}

Two splits are used. The temporal split orders each client's rows by timestamp
within each class and takes the earliest seventy-five per cent for training.
Splitting within class rather than within client is necessary because attacks
occupy contiguous time windows: a flat split by time assigns the test set
whichever attack ran last and places the others entirely in training. The random
split is stratified and serves as the comparison.

Both were measured. The difference between them is 0.0001 macro F1, so the
random split is used for the main results and the temporal split is reported as
an ablation.

\subsection{Evaluation on unseen classes}

Under the cross-test protocol, held-out rows are redistributed across all
clients regardless of what each trained on, so every client is evaluated on
every class. This follows Stones From Other Hills, where a gateway trained on
one set of classes is tested on a partly disjoint set to measure adaptation to
unseen traffic.

The protocol imposes a structural ceiling. A client can only predict classes it
has trained on, so a class held by $h$ of $N$ clients has pooled recall bounded
by $h/N$ and F1 bounded by $2(h/N)/(1 + h/N)$ even under perfect
classification. This bound is computed for every configuration and reported
alongside the achieved value, because a score of 0.34 against a ceiling of 0.41
is a different result from 0.34 against a ceiling of 1.0.

\section{Experimental design}
\label{sec:experimental-design}

Comparisons are paired: both configurations under test see the same partition
and the same initialisation for a given seed, so the difference isolates the
change. Five seeds are used for headline comparisons and three for exploratory
sweeps, with a paired $t$ test. Three seeds is treated as indicative only; at
$n=3$ the critical value is 4.303 and the statistic is dominated by seed
variance.

Results that fall below a difference of 0.03 macro F1 are treated as noise. That
threshold comes from measurement: two runs of an identical configuration
differed by 0.0116, so anything smaller cannot be distinguished from run-to-run
variation.
